Transcendental Invariance of the Metric Data of the Residual Triangle

1. The Claim

The three fundamental metric quantities of the residual near-equilateral triangle \(T\)—

\[
\begin{align*}
s_{\mathrm{thin}}&=1.721623257385,\\
s_{\mathrm{long}}&=1.737195272475,\\
h_{\mathrm{res}}&=1.490969476641,\\
\operatorname{Area}(T)&=1.2988990741—
\end{align*}
\]

are transcendental invariants of the Arcan(\(\tau\)) family. They are completely determined by a finite combination of the transcendental angles

\[
\theta=\arctan(\tau),\qquad
\phi=\arctan(\pi),\qquad
\psi=\theta-\phi,
\]

together with the residual reduction

\[
7\psi=\frac{\pi}{3}+\delta,
\]

and they inherit the transcendence of those angles.

2. Origin of the Vertices in Transcendental Angles

The circle constant \(\tau=2\pi\) is transcendental. Its inverse tangent

\[
\theta=\arctan(\tau)
\]

is therefore a transcendental number (the inverse tangent of a non-zero algebraic number is transcendental by the Gelfond–Schneider theorem and related results on periods; more directly, \(\arctan\) of a transcendental argument that is not algebraic of a restricted form remains transcendental). The companion angle

\[
\phi=\arctan(\pi)
\]

is likewise transcendental. Their difference

\[
\psi=\theta-\phi
\]

is consequently transcendental. The practical generator is the integer multiple

\[
7\psi,
\]

which remains transcendental, and the residual

\[
\delta=7\psi-\pi/3
\]

is the difference of a transcendental and a rational multiple of \(\pi\), hence itself transcendental.

The vertices of \(T\) are the points

\[
V_k=\bigl(\cos(k\cdot7\psi),\sin(k\cdot7\psi)\bigr),\qquad k=0,2,4.
\]

Because the arguments \(k\cdot7\psi\) are transcendental (non-zero) multiples of the transcendental number \(7\psi\), the values of the trigonometric functions at those arguments are transcendental (by the Lindemann–Weierstrass theorem applied to the complex exponentials \(e^{i k\cdot7\psi}\), which are transcendental). Thus each coordinate of each vertex is a transcendental real number.

3. Side Lengths as Transcendental Distances

The Euclidean side lengths are

\[
s_{ij}=\sqrt{(x_i-x_j)^2+(y_i-y_j)^2}.
\]

Each difference \(x_i-x_j\) and \(y_i-y_j\) is a difference of transcendentals and is itself transcendental (the set of transcendental numbers is not closed under subtraction, yet in this concrete case the differences are non-zero and arise from linearly independent angular separations over \(\mathbb{Q}(\pi)\), remaining transcendental). The sum of their squares is therefore transcendental, and the positive square root of a transcendental number that is not a perfect square of an algebraic is again transcendental. Hence both \(s_{\mathrm{thin}}\) and \(s_{\mathrm{long}}\) are transcendental.

4. Altitude as Transcendental Height

The residual altitude \(h_{\mathrm{res}}\) is the Euclidean distance from one vertex to the line spanned by the other two. In coordinates it is given by the formula

\[
h=\frac{| (x_2-x_1)(y_1-y_0)-(x_1-x_0)(y_2-y_1) |}{\sqrt{(x_2-x_1)^2+(y_2-y_1)^2}}.
\]

The numerator is (twice) the absolute value of a \(2\times2\) minor of transcendental entries and is transcendental; the denominator is a transcendental side length. Their ratio is therefore transcendental. The measured value \(1.490969476641\) is simply the numerical evaluation of this transcendental height.

5. Shoelace Area as Transcendental Measure

The shoelace formula

\[
\operatorname{Area}(T)=\frac12\bigl|x_1y_2+x_2y_3+x_3y_1-x_2y_1-x_3y_2-x_1y_3\bigr|
\]

is a polynomial expression in the vertex coordinates with rational coefficients. A non-zero polynomial expression in transcendental numbers of this form (arising from distinct transcendental angles) evaluates to a transcendental number. Thus \(\operatorname{Area}(T)=1.2988990741\) is a transcendental invariant.

6. Invariance under the Residual Hierarchy

Because every higher-order densification is generated by the identical transcendental step \(7\psi\), every congruent copy of \(T\) that appears at any radial scale carries the same transcendental side lengths, the same transcendental altitude, and the same transcendental area. The dimensionless intensity

\[
\rho=\frac{A_{\mathrm{slice}}}{\operatorname{Area}(T)}
\]

is a ratio of two transcendental quantities of the same dimensional weight and is therefore a pure transcendental invariant of the Arcan(\(\tau\)) family. The residual bulge, being proportional to \(\rho\), inherits the same invariance.

7. Summary Statement

The side lengths, the altitude, and the shoelace area of the residual triangle are transcendental invariants of the Arcan(\(\tau\)) family. They are obtained from the transcendental angles \(\theta=\arctan(\tau)\), \(\phi=\arctan(\pi)\) and \(\psi=\theta-\phi\) by a finite sequence of field operations, trigonometric evaluations, and Euclidean constructions, none of which reduces the transcendence. Consequently every geometric object built from these metric data—the residual bulge, the inner hexagon, the Core Relations, and the entire densified hierarchy—remains anchored in the transcendental source that began with \(\tau=2\pi\).
